% Edited conversion of source pp. 18--19.
% The plot is extracted from the source PDF and reconstructed from a base64 asset by LuaLaTeX.
\directlua{
  dofile("scripts/base64_asset.lua")
  decode_base64_file(
    "figures/05-curvature-shells/deflection-curves.jpg.b64",
    "figures/05-curvature-shells/deflection-curves.generated.jpg"
  )
}

\ReportHeading
  {Аналіз впливу кривизни на напружено-деформований стан ортотропних оболонок непрямокутної форми в плані}
  {О.~Я. Григоренко, Є.~В. Філь}
  {Інститут механіки ім. С.~П. Тимошенка НАН України}
  {}

Аналіз напружено-деформованого стану пологих ортотропних оболонок є важливою задачею інженерної практики. Класичні методи здебільшого застосовують до оболонок простої геометрії, що обмежує сферу їх використання. У цій роботі запропоновано розрахункову схему на основі дискретно-континуального підходу, яка дає змогу аналізувати пологі ортотропні оболонки довільної опуклої чотирикутної форми в плані.

Математичну модель побудовано на рівняннях уточненої теорії оболонок, що ґрунтується на гіпотезах прямолінійного елемента (теорія оболонок типу Тимошенка) і враховує поперечні зсуви. Основні рівняння рівноваги та фізичні співвідношення для ортотропного матеріалу зведено до матричної форми
\begin{equation}
  \mathbf{S}\boldsymbol{f}=\boldsymbol{q},
\end{equation}
де вектор-стовпець $\boldsymbol{f}$ містить невідомі функції переміщень $u$, $v$, $w$ та повні кути повороту нормалі $\psi_x$, $\psi_y$.

Ключовим етапом методики є перехід від фізичної області --- довільного опуклого чотирикутника в декартовій системі координат $(x,y)$ --- до канонічної обчислювальної області, одиничного квадрата в локальній системі координат $(\xi,\eta)$. Таке перетворення уніфікує розрахункову схему та дає змогу коректно задавати різноманітні граничні умови на краях оболонок непрямокутної форми в плані.

Двовимірну крайову задачу зводять до одновимірної за допомогою сплайн-апроксимації невідомих функцій за однією з координат, наприклад $\eta$. Невідомі функції подають як лінійні комбінації кубічних B-сплайнів. У результаті вихідна система рівнянь у частинних похідних перетворюється на крайову задачу для системи звичайних диференціальних рівнянь високого порядку відносно координати $\xi$. Отриману систему розв’язують стійким чисельним методом дискретної ортогоналізації.

На основі математичної моделі та чисельного алгоритму створено авторський програмний комплекс. Для верифікації програмного коду проведено тестовий розрахунок напружено-деформованого стану пластини. За нульових головних кривизн, $k_1=k_2=0$, модель пологої оболонки повністю переходить у модель пластини, а результати збігаються з відомими еталонними розв’язками теорії пластин. Це підтвердило коректність матриці диференціальних операторів і врахування граничних умов.

Далі досліджено вплив кривизни на напружено-деформований стан конструкції. Для тієї самої розрахункової області задано $k_1=k_2=0.1$. Наявність кривизни підвищує загальну жорсткість конструкції, тому за однакових зовнішнього навантаження та умов закріплення максимальний прогин пологої оболонки є меншим, ніж прогин плоского аналога.

\begin{center}
  \includegraphics[width=0.58\textwidth]{figures/05-curvature-shells/deflection-curves.generated.jpg}\\[-0.2em]
  {\small Рис.~1. Розподіл безрозмірного прогину $w^{*}=wE/q_0$: 1 --- пластина, $k_1=k_2=0$; 2 --- полога оболонка, $k_1=k_2=0.1$.}
\end{center}

Отже, розроблений підхід і створений програмний комплекс є надійним інструментом для інженерного аналізу оболонок складної геометрії. Вони дають змогу коректно враховувати вплив кривизни, ортотропії матеріалу та форми області в плані на напружено-деформований стан конструкції.

\begin{thebibliography}{9}
\bibitem{curvature:grigorenko2016}
A.~Ya. Grigorenko, W.~H. Müller, Ya.~M. Grigorenko, and G.~G. Vlaikov,
\emph{Recent Developments in Anisotropic Heterogeneous Shell Theory: General Theory and Applications of Classical Theory}.
Berlin: Springer, 2016, 116~p.

\bibitem{curvature:grigorenko2025}
О.~Я. Григоренко, Є.~В. Філь,
``Побудова розрахункової схеми для чисельного аналізу ортотропної непрямокутної в плані пологої оболонки на основі дискретно-континуального підходу,''
у матеріалах конференції молодих учених ``Підстригачівські читання --- 2025,'' Львів, 2025.
\end{thebibliography}
